\section{System Architecture and Authority Topology}\label{sec:design}
\noindent\textbf{Reference Monolithic Workflow.}
A browser client sends a spoken answer to a backend, which transcribes it and calls an \emph{InterviewOrchestrator} (the hub; Fig.~\ref{fig:arch}). The hub owns one session-state dictionary and one lock and runs a fixed per-turn pipeline: (1) a learned \emph{evaluator} (sentence-embedding similarity $S_1$, a concept index $S_2$ and a fine-tuned cross-encoder $R$; $0.15S_1{+}0.35S_2^{\mathrm{eff}}{+}0.50R$ with a cap when a mandatory concept is missed) returns a score and concept evidence; (2) a strict hand-off contract treats a non-finite or out-of-range score as an evaluator failure and fails closed; (3) feedback is produced by a template or an optional local LLM; (4) the hub commits the turn; (5) a baseline schedule and then a frozen PPO policy~\cite{sb3,ppo} propose an easier/same/harder action, followed by a rule layer of six guardrails; (6) a threshold rule decides whether a follow-up question is injected (at most two in a row). Before this work only the PPO policy was a reactive learned decision-maker; the LLM never chose an action. All of this is unchanged when the multi-agent flag is off, and a test shows that with the flag on but no agent contributing the trajectory is identical (Sect.~\ref{sec:results}).

\noindent\textbf{Specialized Agent Roles and Private State.}
\begin{table}[tbp]
\centering\small
\caption{Authority matrix, generated from the enforcing code (\texttt{agents/multiagent/} \texttt{authority.py}). Neither agent has state-write authority; the hub commits alone.}\label{tab:authority}
\begin{tabular}{>{\raggedright\arraybackslash}p{1.5cm}>{\raggedright\arraybackslash}p{2.5cm}>{\raggedright\arraybackslash}p{2.3cm}>{\raggedright\arraybackslash}p{2.8cm}>{\raggedright\arraybackslash}p{2.0cm}}
\toprule
Agent & May propose & Private state & Cannot & On invalid output \\
\midrule
\textbf{Strategy} & difficulty\_action (Easier\slash Same\slash Harder); probe request (concept, kind); strategy mode & weakness ledger (per-concept miss counts, probe history, trend) & write or overwrite any score; set a difficulty level directly; terminate the session; edit history; override hard constraints or guardrail overrides; call another agent & invalid / stale / low-confidence / timed-out proposal is dropped; PPO proposal stands \\
\midrule
\textbf{Coaching} & feedback depth; hint / example flags; concepts to emphasise; remediation tag; tone & coaching history (hints/examples given, repeated misconceptions) & alter evaluator results or scores; cite concepts absent from the evaluator evidence; request difficulty changes; request follow-ups; terminate the session & invalid plan is dropped; the existing template feedback is used unchanged \\
\bottomrule
\end{tabular}

\end{table}
Table~\ref{tab:authority} lists what each agent may propose and may not do. Applying the principle of least privilege~\cite{ss75}, neither agent possesses state-write privileges. Both agents implement a single method \texttt{propose(observation, peer)}, receive an immutable, validated \texttt{TurnObservation} (score, evidence lists, difficulty, score history, progress, follow-up counters; counters carry the value the hub will commit), and return a typed proposal. They are rule- and ledger-based, not LLM-driven; they are the smallest design that has private persistent state, pro-active proposals and a communication channel, which are the properties of weak agency~\cite{wj95} that the existing roles lacked. We make no claim of autonomy beyond that.

\paragraph{Strategy Agent.} Objective: maintain configured difficulty adaptation while tracking persistent weaknesses within the follow-up budget. Private state: per concept, the distinct turns in which the evaluator listed it as missing or incorrect, which concepts were later covered, probes already requested, and an exponential moving average (EMA) of scores. Rules, first match wins: \emph{recover} (current score $<0.35$ and mean of the last two $<0.5$, or a weakness that survived its probes while the score stays $<0.5$) $\rightarrow$ Easier; \emph{remediate} (a concept missed on $\ge 2$ distinct turns and unresolved) $\rightarrow$ Same plus a probe request for that concept; \emph{stretch} (score $\ge0.8$, average $\ge0.7$, last three $\ge0.6$, no open weakness) $\rightarrow$ Harder; otherwise \emph{steady}. Non-\emph{steady} proposals use confidence $\min(0.95, 0.5{+}0.1n)$ for $n$ cited evidence items (a documented heuristic rather than a calibrated probability), while \emph{steady} proposals are assigned 0.45 (below the arbitration threshold $\tau = 0.6$, so the PPO action stands unless the ledger supports a change).

\paragraph{Coaching Agent.} Objective: choose how to coach after an evaluated answer. Private state: concept miss counts and the hints and examples already given. Decisions: up to two concepts to emphasize (incorrect claims first, then repeatedly missed, then evaluator order); a hint if concepts are missing, the score is $<0.6$ and fewer than two hints were given for the lead concept; an example when the lead concept was missed before and no example was given yet; depth (\emph{detailed} if score $<0.4$ or a repeated miss, \emph{brief} if $\ge0.85$ with nothing missing); tone; a remediation tag. A Strategy proposal delivered as a peer message modifies the plan (\emph{recover}: detailed and encouraging; \emph{stretch}: brief; \emph{remediate}: remediation tag on the probed concept). A deterministic renderer applies the plan to the existing feedback dictionary: it adds presentation fields and may prepend a priority tip to, or truncate, the existing improvement and communication-tip lists, and it raises if any protected evaluator field (score, grade, concept and claim lists, transcript, decision source) would change.

\noindent\textbf{Communication and State Ownership.}\label{sec:comm}
Every observation, proposal and decision is a typed object (\texttt{pydantic}, unknown fields rejected, NaN and out-of-range values rejected). The coordinator records each call as an envelope with session, turn, sender, recipient, state version and provenance, and every accept/reject/modify outcome as a decision record with a reason code. Optimistic concurrency uses a \emph{state version} that counts committed turns: a proposal whose version, turn or session differs from the observation, or from the hub's current version, is stale and dropped; a second answer from the same agent for the same version is a duplicate and dropped. The hub builds observations, commits state, increments the version, and persists agent memory in its session snapshot so that a recovered session resumes with the agents' ledgers.

\begin{figure}[tbp]
\centering
\resizebox{0.70\textwidth}{!}{\begin{tikzpicture}[font=\footnotesize,
  box/.style={draw, rounded corners=1.5pt, minimum height=8mm, align=center, inner sep=2pt},
  det/.style={box, fill=gray!12},
  lrn/.style={box, fill=blue!12},
  agt/.style={box, fill=orange!25, very thick},
  arr/.style={-{Stealth[length=2mm]}, thin},
  darr/.style={-{Stealth[length=2mm]}, thin, dashed}]
% outside the hub
\node[det, minimum width=26mm] (ui) at (0.5,6.2) {Browser UI\\(unchanged)};
\node[det, minimum width=26mm] (be) at (3.9,6.2) {Backend\\REST / WebSocket};
\node[lrn, minimum width=30mm] (ev) at (7.5,6.2) {Evaluator\\(learned scorer)};
\node[det, minimum width=30mm] (val) at (11.4,6.2) {Validator +\\strict contract};
\draw[arr] (ui) -- (be); \draw[arr] (be) -- (ev); \draw[arr] (ev) -- (val);
% hub
\draw[rounded corners=3pt, thick, fill=yellow!6] (-0.9,-0.4) rectangle (13.3,5.0);
\node[anchor=north west, font=\footnotesize\bfseries] at (-0.8,4.95) {InterviewOrchestrator (hub: sole writer of session state)};
\node[agt, minimum width=30mm] (st) at (1.6,3.6) {Strategy Agent\\{\scriptsize ledger}};
\node[agt, minimum width=30mm] (co) at (6.9,3.6) {Coaching Agent\\{\scriptsize history}};
\node[lrn, minimum width=30mm] (pp) at (11.2,3.6) {PPO policy\\{\scriptsize (frozen)}};
\node[det, minimum width=71mm] (cd) at (3.85,1.9) {AgentCoordinator: validation $\cdot$ trace};
\node[det, minimum width=44mm] (ar) at (10.1,1.9) {Deterministic Arbitrator\\{\scriptsize Strategy $\leftrightarrow$ PPO}\\{\scriptsize confidence gate $\tau = 0.60$}};
\node[det, minimum width=40mm] (gr) at (10.9,0.3) {Safety Guardrails: final veto\\{\scriptsize $1 \le d \le 5$, $|\Delta d| \le 1$}};
\node[det, minimum width=44mm] (fb) at (2.2,0.3) {Feedback and follow-up\\generators {\scriptsize (LLM optional)}};
\node[det, minimum width=28mm] (cm) at (6.4,0.3) {\textbf{Authoritative}\\\textbf{commit}};
\node[det, minimum width=26mm] (db) at (6.4,-1.4) {SQLite\\persistence};
\draw[arr] (val.south) -- node[right]{\scriptsize committed score} (11.4,5.0);
\draw[darr] (st) -- node[above]{\scriptsize peer message} (co);
\draw[darr] (st.south) -- node[left]{\scriptsize proposal} (1.6,2.35);
\draw[darr] (co.south) -- node[right]{\scriptsize plan} (6.9,2.35);
\draw[arr] (pp.south) -- node[right]{\scriptsize action} (pp.south |- ar.north);
\draw[arr] (cd.east) -- (ar.west);
\draw[arr] (ar.south -| gr.north) -- node[right]{\scriptsize candidate} (gr.north);
\draw[arr] (gr.west) -- (cm.east);
\draw[arr] (2.2,1.45) -- node[right]{\scriptsize plan, probe target} (fb.north);
\draw[arr] (cm.south) -- node[right]{\scriptsize commit} (db.north);
\end{tikzpicture}
}
\caption{System architecture and authority topology of the bounded multi-agent configuration. The \texttt{InterviewOrchestrator} retains sole authority over committed session state. Strategy and Coaching maintain separate private ledgers and contribute proposal-only outputs. Strategy may transmit a unidirectional operational mode to Coaching, while both components submit typed proposals to the coordinator. The PPO policy independently provides the baseline difficulty action. The coordinator routes eligible proposals to deterministic arbitration, after which downstream guardrails retain final veto authority before the hub commits the resulting state transition.}\label{fig:arch}
\end{figure}

\section{Orchestration, Arbitration and Validation}\label{sec:arbitration}
\begin{figure}[tbp]
\centering
\resizebox{0.68\textwidth}{!}{\begin{tikzpicture}[font=\footnotesize, x=1cm, y=1cm,
  life/.style={draw, rounded corners=1pt, minimum width=17mm, minimum height=6mm, align=center, fill=gray!10},
  msg/.style={-{Stealth[length=2mm]}, thin},
  dmsg/.style={-{Stealth[length=2mm]}, thin, dashed},
  note/.style={draw, fill=white, rounded corners=1pt, inner sep=2pt, font=\scriptsize}]
\def\xh{0} \def\xe{2.4} \def\xs{4.8} \def\xc{7.2} \def\xd{9.6} \def\xp{12}
\node[life] (hub) at (\xh,0) {Hub};
\node[life, fill=blue!12] (ev) at (\xe,0) {Evaluator};
\node[life, fill=orange!25] (st) at (\xs,0) {Strategy};
\node[life, fill=orange!25] (co) at (\xc,0) {Coaching};
\node[life] (cd) at (\xd,0) {Coordinator};
\node[life, fill=blue!12] (pp) at (\xp,0) {PPO};
\foreach \x in {\xh,\xe,\xs,\xc,\xd,\xp} \draw[gray, dotted] (\x,-0.4) -- (\x,-8.1);
\draw[msg] (\xh,-0.9) -- node[above]{answer} (\xe,-0.9);
\draw[msg] (\xe,-1.4) -- node[above]{payload} (\xh,-1.4);
\draw[msg] (\xh,-2.0) -- node[above]{1 observation $(v)$} (\xs,-2.0);
\draw[dmsg] (\xs,-2.5) -- node[above]{2 proposal $(v)$} (\xd,-2.5);
\draw[msg] (\xh,-3.1) -- node[above]{3 observation $(v)$ + peer proposal} (\xc,-3.1);
\draw[dmsg] (\xc,-3.6) -- node[above]{4 plan $(v)$} (\xd,-3.6);
\draw[msg] (\xd,-4.1) -- node[above]{validated proposals} (\xh,-4.1);
\node[note] at (\xh+1.35,-4.75) {5 render coaching plan};
\node[note] at (\xh+1.05,-5.3) {6 commit; $v \leftarrow v{+}1$};
\draw[msg] (\xh,-5.9) -- node[above]{7 state observation} (\xp,-5.9);
\draw[msg] (\xp,-6.4) -- node[above]{PPO action} (\xd,-6.4);
\draw[msg] (\xd,-6.9) -- node[above]{8 arbitrated action + guardrail} (\xh,-6.9);
\node[note] at (\xh+1.7,-7.55) {9 follow-up decision, next question};
\end{tikzpicture}
}
\caption{One interview turn as implemented (verbal path). Steps 1--4 happen before the hub commits; arbitration (8) follows the PPO suggestion. On the code-submission path only the Strategy proposal is collected.}\label{fig:flow}
\end{figure}
Figure~\ref{fig:flow} shows one turn. \textbf{Difficulty precedence} (deterministic code, evaluated by ablation rather than assumed optimal): (1) hard constraints (infrastructure failure forces Same; levels stay in $[1,5]$); (2) the guardrail layer, applied last to whichever candidate was selected, wins whenever a rule fires (G3 and G6 hold or raise difficulty, so this layer does not act solely in a downward direction); (3) a validated Strategy proposal with confidence $c \ge \tau$ and cited evidence supersedes the PPO action, recorded as a \emph{conflict}; (4) the PPO action, or the labeled non-RL heuristic when PPO is unavailable. A proposal that cannot take effect at the bounds is recorded as \emph{modified}. \textbf{Probe precedence:} the hard cap of two consecutive follow-ups, then a Strategy probe whose concept appeared in the evaluator's evidence earlier in this session and whose confidence satisfies $c \ge \tau$ (it may add a probe the thresholds would not trigger, or retarget the default one), then the hub's original threshold rule. The retargeted concept is passed to the template follow-up generator; the LLM prompt is not changed.

\textbf{PPO integration.} PPO keeps its role and its 6-D state and 3-action definitions; the reference policy was trained once with a fixed configuration and seed (seed 123, 100,000 nominal training steps) and then frozen before the multi-agent evaluation; no hyperparameter search or result-based checkpoint selection was performed. Its action is one input to arbitration. The agents never see the PPO action before they propose, and the paper makes no claim that the agents improve PPO. Calling this multi-agent RL would be wrong: one policy is learned, and the agents are not trained.

\textbf{Guardrails.} The deployed hub uses an inline copy of the rule layer that differs from the module used by the training environment. A configuration key selects \texttt{legacy\_inline} (default; deployed behavior) or \texttt{canonical}; both are reported separately. On a synthetic grid (39,930 points, hesitation $\mathit{hes}=1-\mathit{conf}$, with $\mathit{conf}$ the acoustic confidence) the two disagree on 4,728 points (11.8\%) in G2, G6 and the no-override case; this is a property of the grid, not a deployment frequency, and a regression test pins it.

\begin{figure}[tbp]
\centering
\resizebox{0.68\textwidth}{!}{\begin{tikzpicture}[font=\scriptsize, node distance=2.2mm,
  s/.style={draw, rounded corners=1pt, minimum height=8mm, minimum width=15.5mm, align=center, inner sep=1.5pt},
  r/.style={draw, dashed, rounded corners=1pt, minimum height=7mm, minimum width=16mm, align=center, fill=red!6, inner sep=1.5pt},
  arr/.style={-{Stealth[length=1.5mm]}, thin}]
\node[s, fill=orange!22] (raw) {agent\\output};
\node[s, right=of raw] (au) {forbidden\\keys?};
\node[s, right=of au] (sc) {schema\\valid?};
\node[s, right=of sc] (sv) {version\\current?};
\node[s, right=of sv] (gr) {concepts\\grounded?};
\node[s, right=of gr, fill=gray!12] (ar) {arbitration\\+ guardrails};
\node[s, right=of ar, fill=green!10] (cm) {hub\\commit};
\foreach \a/\b in {raw/au,au/sc,sc/sv,sv/gr,gr/ar,ar/cm} \draw[arr] (\a) -- (\b);
\node[r, below=6mm of raw] (f0) {timeout,\\exception};
\node[r, below=6mm of au] (f1) {authority\\violation};
\node[r, below=6mm of sc] (f2) {schema\\violation};
\node[r, below=6mm of sv] (f3) {stale\\state};
\node[r, below=6mm of gr] (f4) {ungrounded\\concept};
\foreach \a/\b in {raw/f0,au/f1,sc/f2,sv/f3,gr/f4} \draw[arr, red!60!black] (\a) -- (\b);
\node[below=4mm of f2, text width=9cm, align=center] {each rejection yields no proposal, so the default policy (PPO action, template feedback) decides; an agent never touches the hub's state};
\end{tikzpicture}
}
\caption{Validation pipeline every proposal passes before arbitration. Each rejection yields no proposal; the default policy decides.}\label{fig:valid}
\end{figure}
\paragraph{Proposal Validation.} Before any proposal reaches arbitration, the coordinator executes four sequential checks (Fig.~\ref{fig:valid}): (1)~strict Pydantic schema validation (\texttt{extra=forbid}, NaN and out-of-range rejection); (2)~recursive scanning against the 18 forbidden state keys; (3)~optimistic concurrency verification matching the hub's current committed turn version; and (4)~concept evidence grounding. If an agent raises an unhandled exception or exceeds the configured timeout (150\,ms), its proposal is dropped and the hub proceeds with the default policy.
